\section{Réversible vs irréversible : ce qui se passe, et calcul de $Q$ et $W$ (diapos 91-93)}

\begin{objectif}
\textbf{Objectif :} montrer qu'une transformation \textbf{irréversible} coûte toujours \textbf{plus de
travail} pour obtenir \textbf{moins de chaleur} qu'une transformation réversible entre les deux mêmes états.
C'est la base de tous les rendements du cours (rien n'est jamais parfait).
\end{objectif}

\textbf{Rappel (transformation fermée monotherme, diapo 91) :} le second principe impose $\Delta W\geq0$.
Si la transformation est réversible, la seule possibilité est $\Delta W=0$ et $\Delta Q=0$ ; si elle est
irréversible, la seule possibilité est $\Delta W>0$ et $\Delta Q<0$ (il faut fournir du travail net et on
perd de la chaleur).

\subsection*{Cas de deux transformations réversibles I et II (diapo 92)}
\begin{demotable}
On observe la transformation totale $I+(-II)$ (aller par I, revenir par II en sens inverse) &
C'est un \textbf{cycle fermé monotherme réversible} (composé de deux morceaux réversibles). \\
$W_{tot}=0=W_I+W_{-II}=0 \Rightarrow \boxed{W_I=W_{II}}$ &
D'après le rappel ci-dessus, un cycle fermé monotherme réversible a $\Delta W=0$. Donc les deux chemins
réversibles I et II demandent exactement le \textbf{même travail}. \\
$W_I+Q_I=W_{II}+Q_{II}=\Delta U \Rightarrow \boxed{Q_I=Q_{II}}$ &
Comme $W_I=W_{II}$, l'égalité $\Delta U=W+Q$ (même $\Delta U$ car même trajet A$\to$B) impose aussi
$Q_I=Q_{II}$. \\
\end{demotable}

\subsection*{Cas où I est irréversible et II reste réversible (diapo 93)}
\begin{demotable}
$W_{tot}>0 \Rightarrow W_I+W_{-II}>0 \Rightarrow \boxed{W_I>W_{II}}$ &
Cette fois le cycle total $I+(-II)$ contient une partie irréversible, donc $\Delta W>0$ (rappel plus haut) :
le travail du chemin irréversible I est \textbf{plus grand} que celui du chemin réversible II. \\
$\boxed{W_{\text{irréversible}}>W_{\text{réversible}}}$ et donc $\boxed{Q_{\text{irréversible}}<Q_{\text{réversible}}}$ &
Puisque $\Delta U$ est le même (mêmes points de départ/arrivée) et que $W+Q=\Delta U$ est fixe, si $W$
augmente alors $Q$ doit diminuer d'autant. \\
\end{demotable}

\begin{retenir}
\textbf{Ce qu'il se passe physiquement :} en irréversible, une partie de l'énergie fournie est «gaspillée» (frottements, pertes de charge, combustion non instantanée...). Il faut donc \textbf{fournir un
travail plus important} pour obtenir, au final, \textbf{une quantité de chaleur plus petite} que si la
même transformation avait pu se faire de façon réversible. C'est exactement ce qui justifie les rendements
$\eta_f$, $\eta_{\text{méca}}$ du cours sur les moteurs, et les $\eta_{isos}$ des compresseurs frigo.
\end{retenir}

\newpage
\section{Cycle de Carnot : courbes utilisées et démonstration du rendement (diapos 103-108)}

\begin{objectif}
\textbf{Objectif :} calculer le meilleur rendement \textbf{possible} pour un moteur thermique fonctionnant
entre deux sources de chaleur -- c'est la limite théorique absolue, jamais atteinte en pratique.
\end{objectif}

\textbf{Les 4 courbes du cycle de Carnot} (les deux seules formes de transformation qui n'engendrent
\textbf{aucune irréversibilité}) :
\begin{itemize}[leftmargin=1.4em]
\item A$\to$B : \textbf{compression adiabatique réversible} ($Q=0$, $W>0$)
\item B$\to$C : \textbf{détente isotherme réversible} ($Q>0$, $W<0$) -- la source chaude $T_1$ fournit $Q_1$
\item C$\to$D : \textbf{détente adiabatique réversible} ($Q=0$, $W<0$)
\item D$\to$A : \textbf{compression isotherme réversible} ($Q<0$, $W>0$) -- la source froide $T_2$ reçoit
$Q_2$
\end{itemize}
\emph{En mots :} on comprime et détend le gaz sans jamais échanger de chaleur pendant les phases rapides
(adiabatiques), et on échange toute la chaleur pendant les phases lentes à température constante
(isothermes) -- c'est ce qui rend le cycle totalement réversible.

\begin{demotable}
$\eta=\dfrac{\text{énergie utile}}{\text{énergie fournie}}=\dfrac{|W_{tot}|}{Q_1}=\dfrac{|-Q_1-Q_2|}{Q_1}$ &
Le rendement, c'est toujours ce qu'on récupère (le travail net) divisé par ce qu'on a dû payer
(la chaleur $Q_1$ fournie par la source chaude). $W_{tot}=-Q_1-Q_2$ vient de $\Delta U_{cycle}=0=Q_1+Q_2+W_{tot}$. \\
$\eta=\dfrac{Q_1+Q_2}{Q_1}=\boxed{1+\dfrac{Q_2}{Q_1}}$ \quad ($W_{tot}<0$ car machine motrice) &
On simplifie la valeur absolue car $W_{tot}$ est négatif pour un cycle moteur. \\
Comme $Q_2<0$ et $|Q_2|<|Q_1|$ : $0\leq\eta<1$ &
$Q_2/Q_1$ est négatif mais de valeur absolue $<1$, donc $\eta$ reste entre 0 et 1 : jamais 100\%. \\
Sur une isotherme, $\Delta U=0 \Rightarrow Q=-W=-nRT\ln\!\left(\dfrac{V_i}{V_f}\right)$ &
Pour calculer $Q_1$ et $Q_2$ explicitement, on utilise le travail d'une isotherme (section 1) et le fait
que $Q=-W$ quand $\Delta U=0$. \\
$\eta=1+\dfrac{Q_{DA}}{Q_{BC}}=1-\dfrac{T_A\ln(V_A/V_D)}{T_C\ln(V_B/V_C)}$ &
On applique la formule ci-dessus aux deux isothermes B-C (à $T_C$) et D-A (à $T_A$) ; les $Q_1,Q_2$
deviennent des logarithmes de rapports de volumes. \\
Relations adiabatiques : $T_AV_A^{\gamma-1}=T_BV_B^{\gamma-1}$ (A-B) et $T_CV_C^{\gamma-1}=T_DV_D^{\gamma-1}$
(C-D), avec $T_B=T_C$ et $T_D=T_A$ &
On utilise la 1\textsuperscript{ère} formule de l'adiabatique (section 2) sur les deux portions adiabatiques
du cycle, en notant que B et C sont à la même température (idem D et A, isothermes). \\
$\left(\dfrac{V_A}{V_B}\right)^{\gamma-1}=\dfrac{T_B}{T_A}=\dfrac{T_C}{T_D}=\left(\dfrac{V_D}{V_C}\right)^{\gamma-1}
\Rightarrow \dfrac{V_A}{V_B}=\dfrac{V_D}{V_C} \Rightarrow \boxed{\dfrac{V_A}{V_D}=\dfrac{V_B}{V_C}}$ &
En combinant les 4 relations, on montre que les deux rapports de volumes sont \textbf{égaux} -- c'est
l'étape clé : cela va faire disparaître les logarithmes de volume dans le rendement ! \\
Donc $\ln(V_A/V_D)=\ln(V_B/V_C)$, ces deux termes se simplifient dans la fraction &
Puisque les rapports sont égaux, leurs logarithmes le sont aussi : la fraction $\frac{\ln(V_A/V_D)}{\ln(V_B/V_C)}$
vaut exactement $1$. \\
$\boxed{\eta_{Carnot}=1-\dfrac{T_A}{T_C}}$ &
Il ne reste que le rapport des températures. Avec $T_A=T_2$ (froide) et $T_C=T_1$ (chaude), c'est la formule
$\eta_{Carnot}=1-T_2/T_1$ du \textbf{formulaire}. \\
\end{demotable}

\begin{retenir}
Même dans ce cas théorique \textbf{parfait} (aucune irréversibilité), le rendement n'est jamais égal à 1 :
il ne dépend que du rapport des deux températures des sources. C'est la limite indépassable pour
\textbf{tout} moteur thermique ditherme, quel que soit le fluide ou la technologie utilisée.
\end{retenir}

\newpage
\section{Cycle de Joule (turbine à gaz) : démonstration et utilisation (diapos 109-116)}

\subsection*{Le cycle fermé théorique (diapos 109-112)}

\textbf{Les 4 transformations} (même schéma logique que Carnot, mais avec des \textbf{isobares} à la place
des isothermes -- plus simple à réaliser techniquement, mais irréversibles) :
\begin{itemize}[leftmargin=1.4em]
\item A$\to$B : compression \textbf{adiabatique réversible} ($Q=0$, $W>0$)
\item B$\to$C : détente \textbf{isobare irréversible} ($Q>0$, $W<0$) -- la source chaude $T_1$ fournit $Q_1$
\item C$\to$D : détente \textbf{adiabatique réversible} ($Q=0$, $W<0$)
\item D$\to$A : compression \textbf{isobare irréversible} ($Q<0$, $W>0$) -- la source froide $T_2$ reçoit
$Q_2$
\end{itemize}

\begin{demotable}
$\eta=1+\dfrac{Q_2}{Q_1}=1+\dfrac{Q_{DA}}{Q_{BC}}=1+\dfrac{nC_p(T_A-T_D)}{nC_p(T_C-T_B)}$ &
Même point de départ que Carnot ($\eta=1+Q_2/Q_1$), mais ici les échanges de chaleur se font à
\textbf{pression constante} : on utilise donc $Q=nC_p\Delta T$ au lieu de la formule isotherme. \\
Relations adiabatiques (A-B et C-D) : $T_Ap_A^{(\gamma-1)/\gamma}=T_Bp_B^{(\gamma-1)/\gamma}$ et
$T_Dp_D^{(\gamma-1)/\gamma}=T_Cp_C^{(\gamma-1)/\gamma}$, avec $p_A=p_D$ et $p_B=p_C$ (isobares) &
On utilise la 3\textsuperscript{ème} formule de l'adiabatique ($T^\gamma p^{1-\gamma}=$cste) sur les deux
portions adiabatiques ; les pressions se correspondent car B-C et D-A sont isobares. \\
$\Rightarrow \dfrac{T_A}{T_D}=\dfrac{T_B}{T_C}=\dfrac{T_A-T_D}{T_B-T_C}=-\dfrac{T_A-T_D}{T_C-T_B}$ &
En combinant les deux relations (même astuce que pour Carnot : deux rapports égaux $\Rightarrow$ leur
différence donne le même rapport), on relie directement les 4 températures. \\
$\boxed{\eta_J=1-\dfrac{T_A}{T_B}}$ &
Les $(T_A-T_D)$ se simplifient avec ceux du rendement, il ne reste que $T_A$ et $T_B$. \\
Comparaison à Carnot : $\eta_{Carnot}=1-\dfrac{T_A}{T_C}$, or $T_C>T_B$ &
$T_C$ est la température atteinte après la détente isobare complète, $T_B$ seulement après la compression :
$T_C>T_B$ toujours. \\
$\boxed{\eta_{Carnot}>\eta_{Joule}}$ &
Puisque $T_C>T_B$, diviser $T_A$ par $T_C$ donne un nombre plus petit que diviser $T_A$ par $T_B$ : le
rendement de Joule est \textbf{plus bas} que celui de Carnot, car les isobares ne sont pas réversibles. \\
\end{demotable}

\subsection*{Utilisation : le cycle de Joule ouvert (groupe électrogène / turboréacteur, diapos 113-116)}

\textbf{En pratique}, on ne referme pas le cycle : un compresseur axial (entraîné par la turbine, sur le
même arbre) aspire de l'air frais en A, l'air ressort comprimé en B, est chauffé dans la chambre de
combustion (injection de carburant, $Q_1$) jusqu'en C, puis se détend dans la turbine jusqu'en D en
ressortant sous forme de gaz brûlés (remplacés par de l'air neuf en A au cycle suivant -- \textbf{cycle
ouvert}). \textbf{Une grande partie du travail récupéré à la détente sert directement à entraîner le
compresseur.}

\begin{demotable}
$\eta=\dfrac{-W_{ut}}{Q_1}=\dfrac{-(W_{CD}-W_{AB})}{Q_{BC}}$, avec $Q_{BC}=H_C-H_B=nC_p(T_C-T_B)$ &
Ici les machines sont \textbf{ouvertes} (compresseur, turbine) : on utilise le rendement basé sur
$W_{ut}$, et $Q_{BC}$ (isobare) se calcule avec l'enthalpie, comme pour tout échangeur ouvert. \\
Machine ouverte adiabatique : $H_f-H_i=W_{ut}+Q$ et $Q=0 \Rightarrow W_{ut}=\Delta H$ &
Résultat de la section 1 (machine ouverte) appliqué aux étapes adiabatiques A-B et C-D : pas de chaleur
échangée, donc tout le travail utile = variation d'enthalpie. \\
$W_{CD}=H_D-H_C$ et $W_{AB}=H_B-H_A$ &
On applique ce résultat spécifiquement aux deux détentes/compressions adiabatiques du cycle. \\
$\eta=\dfrac{-H_D+H_C-H_B+H_A}{H_C-H_B}$, et $\Delta H=nC_p\Delta T$ &
On remplace $W_{CD}$ et $W_{AB}$ par leurs expressions en $H$, puis on repasse en températures. \\
$\boxed{\eta=1+\dfrac{-T_D+T_A}{T_C-T_B}}$ &
Même résultat que la version fermée (à une réécriture près) : la démonstration par les enthalpies confirme
le résultat obtenu avec les températures. \\
\end{demotable}
